% Part: many-valued-logic
% Chapter: syntax-and-semantics
% Section: semantic-notions

\documentclass[../../../include/open-logic-section]{subfiles}

\begin{document}

\olfileid[hi]{mvl}{syn}{sem}

\olsection{अर्थवैज्ञानिक धारणाएँ}

मान लें बहुमूल्य तर्कशास्त्र $\Log L$ किसी आव्यूह द्वारा दिया गया है। तब
हम~$\Log L$ के लिए प्रचलित अर्थवैज्ञानिक धारणाएँ परिभाषित कर सकते हैं।

\begin{defn} 
\begin{enumerate}
\item !!^a{formula}~$!A$ \emph{संतुष्ट्य} है यदि किसी~$\pAssign{v}$ के लिए
  $\pSat{v}{!A}$; और \emph{असंतुष्ट्य} है यदि किसी भी $\pAssign{v}$ के लिए
  $\pSat{v}{!A}$ न हो;
\item !!^a{formula}~$!A$ \emph{सर्वसत्यता} है यदि सभी
  !!{valuation}~$v$ के लिए $\pSat{v}{!A}$;
\item यदि $\Gamma$ !!{formula} का कोई समुच्चय है, तो $\Gamma \Entails !A$
  (``$\Gamma$ से $!A$ का अनुगमन होता है'') तभी और केवल तभी जब प्रत्येक
  ऐसे !!{valuation}~$\pAssign{v}$ के लिए $\pSat{v}{!A}$ जिसके लिए
  $\pSat{v}{\Gamma}$।
\item यदि $\Gamma$ !!{formula} का कोई समुच्चय है, तो $\Gamma$
  \emph{संतुष्ट्य} है यदि कोई ऐसा !!a{valuation}~$\pAssign{v}$ हो जिसके
  लिए $\pSat{v}{\Gamma}$; अन्यथा $\Gamma$ \emph{असंतुष्ट्य} है।
\end{enumerate} 
\end{defn}

इन धारणाओं के लिए हमारे पास चिरसम्मत तर्कशास्त्र वाली स्थिति के कुछ समान
तथ्य हैं:

\begin{prop}
\ollabel{prop:semanticalfacts} 
\begin{enumerate} 
\item $!A$ सर्वसत्यता तभी और केवल तभी है जब $\emptyset \Entails !A$;
\item यदि $\Gamma$ संतुष्ट्य है, तो $\Gamma$ का प्रत्येक परिमित उपसमुच्चय
  भी संतुष्ट्य है;
\item\ollabel{def:monotonicity}%
एकदिष्टता: यदि $\Gamma \subseteq \Delta$ और $\Gamma \Entails !A$, तो
  $\Delta \Entails !A$ भी;
\item\ollabel{def:Cut}%
संक्रामकता: यदि $\Gamma \Entails !A$ और $\Delta \cup \{ !A\} \Entails
  !B$, तो $\Gamma \cup \Delta \Entails !B$;
\end{enumerate}
\end{prop}

\begin{proof}
अभ्यास।
\end{proof}

\begin{prob}
\olref[mvl][syn][sem]{prop:semanticalfacts} सिद्ध करें।
\end{prob}

चिरसम्मत तर्कशास्त्र में हम अनुगमन और सप्रतिबंध को जोड़ सकते हैं।
उदाहरणार्थ, \emph{मोदुस पोनेन्स} की वैधता है: यदि $\Gamma \Entails !A$ और
$\Gamma \Entails !A \lif !B$, तो $\Gamma \Entails !B$। चिरसम्मत
तर्कशास्त्र में $\Entails$ और $\lif$ का दूसरा महत्त्वपूर्ण संबंध
अर्थवैज्ञानिक निगमन प्रमेय है: $\Gamma \Entails !A \lif !B$ तभी और केवल
तभी जब $\Gamma \cup \{!A\} \Entails !B$। ये परिणाम बहुमूल्य तर्कशास्त्रों
में सदा लागू \emph{नहीं} होते। उनका लागू होना सत्य-फलन~$\tf{\lif}$ पर निर्भर
है।

\end{document}
